%% Generated from paper/draft/03_mechanisms.md by md2tex.py.
%% Mechanical first pass -- polish by hand, then stop regenerating.

\section{Three mechanisms of the return-arrival process}
\label{sec:03}

\subsection{Grid alignment decides the \emph{sign} of the gap effect}

Comparing an agentic arm (turns separated by tool gaps) against a conventional arm (same work, no gaps) across concurrency levels does not give a single direction. At some concurrencies the gap makes utilization \emph{better}, at others worse. The distinguishing variable is not load but \textbf{where the offered concurrency falls on the compiled grid}: when it lands between two rungs, the gap breaks the running set into pieces that pad less, and the arm with gaps wins.
% CLAIMS 1.5

We establish this causally. Holding the workload, the seed, the model and the slot count fixed, we recompiled the artifact with one additional rung and measured again. On the grid \texttt{(1,2,4,8)} the pooled ratio at N=6 was 1.1504; adding rung 6 to make \texttt{(1,2,4,6,8)} moved it to 0.9717, and the reversal disappeared. An ablation completes the attribution: on a hypothetical continuous grid the law vanishes to exactly 1.0000, since with no padding there is no padding difference to explain (model).
% CLAIMS 1.6
% CLAIMS 1.7

\begin{figure}[t]
  \centering
  \includegraphics[width=\columnwidth]{fig2_grid_alignment}
  \caption{Whether a tool gap helps or hurts is decided by where the offered concurrency falls on the compiled grid, and the effect changes sign. Pooled utilization ratio (agentic/conventional) at an admission ceiling of eight; the square series is the same workload, seed, model and slot count after a recompile that added one rung, which is what makes the attribution causal rather than correlational. The flat line at 1.0000 is an ablation on a continuous grid (model).}
  \label{fig:fig2}
\end{figure}


\subsection{The reuse cliff counts requests, not tokens}

A completed prefix either survives a gap intact or is gone; there is no partial survival, because a request of 8,192 tokens or fewer occupies exactly one outer slot regardless of its length. Survival is therefore decided by \textbf{how many other requests were admitted during the gap}, not by how many tokens they carried. On this artifact --- eight outer slots of 8,192 tokens under FIFO, with 2,000-token requests --- reuse survives six background requests and dies at the seventh, and that threshold reproduced identically in 12 of 12 trials.
% CLAIMS 1.1
% CLAIMS 1.2

The eviction order is FIFO over \emph{allocation}, not over last use, so the oldest allocation is the first victim and re-access does not protect it. An ablation replacing the policy with LRU over blocks removes the count threshold entirely and makes the threshold inversely proportional to request size instead (model).
% CLAIMS 1.3

One asymmetry is worth stating because it makes gaps costlier than a naive account suggests: a session that returns \emph{immediately} does not evict its own cache, because a completed block becomes evictable only after the next admission has already chosen its victim. A session that returns after a gap has lost that protection. \textbf{The cost of a gap therefore includes losing one's own shield, not merely gaining competitors.}
% CLAIMS 1.13

\begin{figure}[t]
  \centering
  \includegraphics[width=\columnwidth]{fig1_survival_cliff}
  \caption{Reuse of a completed prefix is all-or-nothing, and the threshold counts requests rather than tokens. Measured on the layer-2 pool of this artifact (eight outer slots of 8{,}192 tokens, FIFO) with 2{,}000-token requests, where reuse survives six background arrivals and disappears at the seventh, reproduced identically in 12 of 12 trials. The grey step is the upper-layer metric, which keeps reporting hits after the tensors are gone; the dashed steps are an ablation (model, not measurement) replacing FIFO with block-granular LRU, under which the threshold becomes size-dependent.}
  \label{fig:fig1}
\end{figure}


\subsection{Prefill is charged to everyone}

Injecting a prefill while several sessions decode produces a latency spike in \textbf{all} of them, beginning within a millisecond of each other and lasting the prefill's duration; across nine runs the ratio of spike to prefill time stayed in 1.01--1.14. The stall is predictable from the computed token count: on this hardware and model, \texttt{ceil(n/128) $\times$ (0.0212 + 6.4e-7$\cdot$n)} seconds fits four injection sizes from 500 to 6,000 tokens with a worst residual of 2.4 ms.
% CLAIMS 1.8
% CLAIMS 1.9

Adding this term is what makes a cost model transfer across concurrency. A decode-only model's predicted/measured ratio ran from 0.565 to 0.855 with a monotone dependence on concurrency; adding the serialisation term moves it to 0.973--1.039 and removes the dependence, accounting for 87--120 \% of the earlier bias.
% CLAIMS 1.10

\begin{figure}[t]
  \centering
  \includegraphics[width=\columnwidth]{fig3_prefill_tax}
  \caption{Prefill is charged to every concurrent session, and a cost model that omits it is biased in a way that grows with concurrency. Predicted over measured inter-token-latency sum for both arms; hollow markers omit the serialisation term and filled markers include it, on this hardware and model.}
  \label{fig:fig3}
\end{figure}


\subsection{Why the three multiply}

The mechanisms are not three independent costs to be summed. The slot pool decides \textbf{who} must recompute; exclusive prefill decides that the recomputation is charged to \textbf{everyone then decoding}; and the grid decides \textbf{how many} that is. A cache miss under a wide batch is expensive twice over.

This composition is also why an intuition that works elsewhere fails here, and Section~\ref{sec:05} takes it up: holding a returning session back so that it can share a batch looks like it should help, and does not.
